\chapter{Introducción}
\label{cap:introduccion}

La protección y la supervisión de espacios interiores son una necesidad constante
en entornos tan diversos como oficinas, naves industriales, almacenes o viviendas.
Este capítulo enmarca el trabajo en ese ámbito: se argumenta la importancia y la
actualidad del problema, se delimita el reto técnico que se aborda y se enuncia el
propósito del presente Trabajo Fin de Grado, para terminar describiendo la
organización de la memoria.

%% =============================================================================
\section{Importancia y actualidad de la temática}
\label{sec:importancia-actualidad}

La vigilancia de espacios interiores cerrados se apoya tradicionalmente en redes
de cámaras fijas y en la supervisión por parte de personal de seguridad. Ambos
recursos presentan limitaciones conocidas: las cámaras fijas solo cubren los
ángulos en los que se han instalado y dejan zonas ciegas, y la atención humana
sostenida sobre múltiples monitores resulta costosa y propensa a la fatiga. Como
consecuencia, una parte de los eventos relevantes puede pasar inadvertida o
detectarse con retraso, justo cuando una respuesta temprana es más valiosa.

El interés por automatizar esta tarea es especialmente vigente hoy en día. En los
últimos años, la robótica móvil ha alcanzado un grado de madurez que
permite delegar parte de la vigilancia en plataformas autónomas: la consolidación
de \gls{ros} como middleware estándar para el desarrollo de aplicaciones
robóticas, con despliegues reales documentados \parencite{macenski2022ros2}, junto
con la aparición de pilas de navegación robustas \parencite{macenski2020nav2} y de
hardware asequible, ha rebajado de forma notable la barrera de entrada. De manera
complementaria, los detectores de objetos en tiempo real basados en redes
neuronales convolucionales \parencite{redmon2016yolo}, hoy disponibles en
implementaciones eficientes y de uso extendido \parencite{jocher2023ultralytics},
hacen viable dotar a un robot de una capacidad de percepción suficiente para
identificar situaciones de interés mientras se desplaza.

La convergencia de ambos avances habilita un modelo de vigilancia activa, en el
que un agente móvil recorre el espacio de forma autónoma y reclama la atención del
operador únicamente ante los eventos que detecta. Frente al esquema pasivo de la
cámara fija, este enfoque amplía la cobertura efectiva con un único dispositivo en
movimiento y filtra la información que llega al usuario. Es en esta ventana de oportunidad 
donde se sitúa el presente trabajo.

%% =============================================================================
\section{Planteamiento del problema}
\label{sec:planteamiento-problema}

Construir un robot de patrulla autónoma para interiores no constituye un problema
único, sino la integración coordinada de varios subproblemas. El robot debe
cartografiar un entorno inicialmente desconocido, planificar una ruta que cubra el
espacio de forma razonable \parencite{galceran2013coverage}, navegar evitando
obstáculos, percibir la presencia de personas y comunicar las incidencias a un
sistema con el que el usuario pueda interactuar. Cada una de estas capacidades
descansa en tecnologías distintas (navegación, percepción, mensajería, servicios
web y aplicaciones de usuario) que rara vez comparten el mismo plano de ejecución.

El reto principal, por tanto, no es resolver cada subproblema de manera aislada,
sino integrarlos en un sistema coherente que sea a la vez funcional, seguro y
mantenible. Un acoplamiento excesivo entre el robot y el servidor, en el que cada extremo
depende de la disponibilidad inmediata del otro, penalizaría la escalabilidad y
la robustez del sistema y elevaría el consumo de recursos de una plataforma con
autonomía energética limitada. De ahí la necesidad
de crear una arquitectura que desacople la capa física del robot del plano de control y
de los servicios de usuario.

%% =============================================================================
\section{Propósito y finalidad del trabajo}
\label{sec:proposito}

El propósito de este Trabajo Fin de Grado es diseñar e implementar un sistema
robótico de patrulla autónoma para estancias cerradas que dé respuesta al problema
anterior, construido sobre un robot \mbox{TurtleBot~4} y \gls{ros}. El sistema
cartografía el entorno mediante \gls{slam}, deriva del mapa resultante una ruta de
cobertura y la recorre con la pila de navegación \gls{nav2}. Durante el recorrido,
un módulo de percepción basado en \gls{yolo} detecta la presencia de personas,
graba un clip del evento y lo registra como incidencia.

La finalidad última no es solo que el robot patrulle, sino demostrar la viabilidad
de una arquitectura distribuida y orientada a eventos que integre la capa
robótica, el procesamiento a bordo y los servicios web en un único sistema
funcional, extensible y seguro. La comunicación entre el robot y el resto del
sistema se resuelve mediante un bus de mensajería \parencite{kreps2011kafka} que
desacopla el hardware del backend, de modo que el robot solo transmite
información cuando se produce un evento relevante. Un servidor
desarrollado con FastAPI consume estos eventos, persiste la información y expone
una \gls{api} \gls{rest} que una aplicación \mbox{Flutter} emplea para que el
usuario controle la patrulla y revise el historial de incidencias. El proyecto se
organiza en cuatro bloques diferenciados (capa robótica,
infraestructura de datos, backend y aplicación de usuario) desplegables mediante
contenedores.

%% =============================================================================
\section{Estructura de la memoria}
\label{sec:estructura-memoria}

El resto del documento se organiza como sigue. El
Capítulo~\ref{cap:marco-teorico} revisa los fundamentos teóricos y el estado del
arte en navegación autónoma, percepción visual y arquitecturas orientadas a
eventos. El Capítulo~\ref{cap:objetivos} formula el objetivo general y los
objetivos específicos del trabajo. El Capítulo~\ref{cap:metodologia} describe la
metodología de desarrollo, las herramientas empleadas y los criterios de
evaluación previstos. El Capítulo~\ref{cap:desarrollo} constituye el núcleo
técnico de la memoria y detalla la arquitectura del sistema y la implementación
de cada uno de sus subsistemas. El Capítulo~\ref{cap:resultados} presenta los
resultados obtenidos y su análisis. Finalmente, el
Capítulo~\ref{cap:conclusiones} recoge las conclusiones, las limitaciones
detectadas y las líneas de trabajo futuro.
